\section{Fraud scenarios and variants}
\label{app:scenarios}

The \nScenarios{} scenarios and their main signals, as the scenario design notes give them
(\texttt{docs/\allowbreak ml/\allowbreak scenarios/}), with the fraud rows each contributes to the draw \texttt{d8083dbc}
(Table~\ref{tab:scenarios}). The notes describe what a detection system can observe, not how to
commit fraud. Of the \nProvFraudTotal{} fraud parameters, \nProvFraudAssumed{} are assumed and
\nProvFraudCalibrated{} are calibrated to specification targets
(\texttt{dataset/\allowbreak params\_provenance.md}); none is taken from a published East African
fraud statistic. Agent, card-not-present and merchant-fraud incidents can be a single transaction;
the other scenarios' incidents always span several (\texttt{fraud.rows\_per\_incident}).

\begin{longtable}{p{0.2\linewidth}p{0.58\linewidth}r}
\toprule
Scenario & Main signals & Fraud rows \\
\midrule
\endhead
SIM-swap takeover & A SIM swap shortly before transfers to new or mule counterparties. & \nScenSimSwap{} \\
Account takeover & A new device, then online spend and transfers. & \nScenATO{} \\
Agent fraud & Cash-outs at a small set of compromised agents, in round sums. & \nScenAgent{} \\
Velocity fraud & Bursts of small payments within minutes. & \nScenVelocity{} \\
Card-not-present & Online or card spend, often from a device the account has never used. & \nScenCNP{} \\
Mule accounts & Many inbound drains and fast forwarding outward (fan-in and fan-out). & \nScenMule{} \\
Merchant fraud & Card charges at a small set of colluding merchants. & \nScenMerchant{} \\
Synthetic identity & A quiet new account, then a bust-out. & \nScenSynth{} \\
\midrule
Total & & \nScenTotal{} \\
\bottomrule
\end{longtable}

Shared by every scenario: fraud rows go through the same code path, formats and null rules as
legitimate rows, and incidents follow the victim's daily activity profile, with a share at night.

\paragraph{Variants.} These variants are scored separately in Section~\ref{sec:generalisation}.
\begin{itemize}
    \item \texttt{adapted\_below\_threshold}: from the \nAdaptMonthOrdinal{} simulated month, some amounts
  adapt to just below one simulated, assumed rule threshold. It is present in training and test.
    \item \texttt{novel\_esim\_delayed\_drain}: a SIM-swap sub-variant whose enabling event is
  recorded as a device change, whose drain follows it by days, and which runs on bank transfer for
  customers with a smartphone; \nNovelRows{} rows, test period only.
  \item \texttt{reversal\_scam\_social\_engineering}: the pre-specified non-burst variant
  (ADR 0028); \nRevRows{} rows, test period only, generated as \texttt{mule\_account} rows,
  paid to a counterparty on the victim's payee list (our modelling choice; ADR 0028's citation of
  the scam typology for it is not supported by the cited sources), with its own amount rule; \nRevNovelScored{} of the \nRevFraud{} scored rows go to a payee the
  account had never paid (Section~\ref{sec:generalisation}). The battery scores \nRevFraud{} of them as fraud: the
  other \nRevScoredUnlabelled{} are fraud whose observed label says legitimate, the benchmark's
  missed-fraud noise (Section~\ref{sec:generalisation}).
\end{itemize}
